\documentclass[11pt,a4paper]{article}

\usepackage[margin=2.5cm]{geometry}
\usepackage{booktabs}
\usepackage{float}
\usepackage{tabularx}
\usepackage{array}
\usepackage{enumitem}
\usepackage{totcount}
\usepackage{fancyhdr}
\usepackage{lastpage}
\usepackage[hidelinks]{hyperref}
\setlength{\emergencystretch}{1.5em}

% Section numbers in the same style as the ConOps: 1.0, 1.1.
\renewcommand{\thesection}{\arabic{section}.0}
\renewcommand{\thesubsection}{\arabic{section}.\arabic{subsection}}

% Controlled-document header on every page.
\newcommand{\docref}{Reference TBD}
\newcommand{\docrev}{Issue 1, Rev.~0 --- DRAFT}
\pagestyle{fancy}
\fancyhf{}
\lhead{\small Remote Controller with Fail-Safe Update --- Software Classification}
\rhead{\small \docrev}
\lfoot{\small \docref}
\rfoot{\small Page \thepage\ of \pageref{LastPage}}

% Open items: values not yet known, listed in App. A; the total is counted.
\newtotcounter{tbd}
\newcommand{\tbditem}[2]{\refstepcounter{tbd}\label{tbd:#1}\item[TBD-\thetbd] #2}
\newcommand{\tbd}[1]{\mbox{\textbf{[\hyperref[tbd:#1]{TBD-\ref*{tbd:#1}}]}}}

\setlist[description]{style=multiline,leftmargin=2cm,nosep,font=\normalfont\bfseries}
\setlist[itemize]{nosep}

\begin{document}

\thispagestyle{fancy}
\begin{center}
{\LARGE\bfseries Remote Controller with Fail-Safe Update\par}
\vspace{0.4em}
{\Large Software Classification\par}
\vspace{1em}
\begin{tabular}{@{}ll@{}}
Reference & TBD \tbd{ref} \\
Issue / revision & Issue 1, Revision 0 --- DRAFT \\
Date & 10 October 2026 \\
Author & Andrew Toma \\
Open items & \total{tbd} TBD (App.~A) \\
\end{tabular}
\end{center}

\subsection*{Document change record}
\begin{tabular}{@{}llll@{}}
\toprule
Issue & Rev. & Date & Change \\
\midrule
1 & 0 & 2026-10-10 & First draft \\
\bottomrule
\end{tabular}

% ===========================================================================
\section{Purpose and Scope}
This document establishes the software classification of each software item,
and so the rigour each item requires. It is the formal record of the software
classification determinations [1, Sec.~3.5.2, p.~25].

The evaluated software items encompass the System software and its enabling
systems, as defined in the Concept of Operations [4, Sec.~1.2.2]. The System
models an uncrewed spacecraft whose primary mission objectives directly depend
on this software; a failure would cost those objectives.

% ===========================================================================
\section{Documents}
\noindent [1] NASA, \emph{NASA Software Engineering Requirements}, NPR 7150.2D.

\noindent [2] NASA, \emph{Software Assurance and Software Safety Standard},
NASA-STD-8739.8B, 2022.

\noindent [3] NASA, \emph{NASA Software Engineering Handbook}, NASA-HDBK-2203,
Version~D, page ``SWE-020 -- Software Classification''.

\noindent [4] \emph{Remote Controller with Fail-Safe Update --- Concept of
Operations}, Issue~1, Revision~0, Draft.

% ===========================================================================
\section{Classification}
Each item takes the highest class whose definition applies
[1, Sec.~3.5.1, p.~24].

\begin{table}[H]
\centering
\begin{tabular}{@{}lll@{}}
\toprule
Software item & Role & Class \\
\midrule
Controller firmware & FLIGHT & B \\
Onboard computer software & FLIGHT & B \\
Ground software in the mission operations center & GROUND & B \\
Bench test tools & TEST & C \\
Firmware build tools & BUILD & D \\
\bottomrule
\end{tabular}
\end{table}

\noindent\textbf{Legend: role}
\begin{description}[leftmargin=2.4cm]
\item[FLIGHT] Operates the spacecraft from orbit.
\item[GROUND] Operates the spacecraft from the ground.
\item[TEST] Tests the System; never operates the spacecraft.
\item[BUILD] Builds the System; never operates the spacecraft.
\item[Class] As defined in [1, App.~D, pp.~79--86].
\end{description}

\subsection{Rationale}
\noindent\textbf{Controller firmware: Class B.} Class B covers flight software
  ``required to operate the vehicle or space asset \ldots{} such as commanding
  of the vehicle or asset'' [1, App.~D, p.~81]; its examples include ``fault
  protection'' [1, App.~D, p.~81]. The controller runs commands and returns to
  working firmware on its own when new firmware fails.\par\smallskip
\noindent\textbf{Onboard computer software: Class B.} The same definition applies. It
  passes every command and firmware image to the controller, watches it, and
  restores it when its firmware cannot run.\par\smallskip
\noindent\textbf{Ground software: Class B.} Class B examples include ``command and control
  ground systems'' [1, App.~D, p.~81]. The ground software prepares and checks
  every command and firmware image sent to the spacecraft.\par\smallskip
\noindent\textbf{Bench test tools: Class C.} Class C covers ``software used for the testing
  of space assets'' [1, App.~D, p.~82], and Class B excludes systems ``used in
  testing Class B systems'' in a development environment [1, App.~D, p.~81].\par\smallskip
\noindent\textbf{Firmware build tools: Class D.} Class D covers ``ground software tools that
  support engineering development'' [1, App.~D, p.~83].\par\smallskip

% ===========================================================================
\section{Safety-Critical Designation}
Some software is safety-critical: the parts that can cause a hazard, or that
exist to prevent, detect or handle one [2, Sec.~4.2, p.~20]. Safety-critical
software must meet extra requirements on top of its class, for example test
coverage of every decision (MC/DC) [1, Sec.~3.7.4, pp.~26--27].

Which items are safety-critical can only be decided once the hazards are known,
in the hazard analysis \tbd{safety}. Every item here is already Class~D or
higher, as safety-critical software must be [3, Sec.~3.2.2], so the result
cannot change a class.

% ===========================================================================
\section{Tailoring and Review}
\begin{itemize}
\item Software engineering and software assurance are expected to agree on each
  classification [1, Sec.~3.5.1 Note, p.~25]. One engineer builds the System,
  with no independent assurance team (ConOps constraint C-3), so the
  classification is the author's alone.
\item The classification is revisited at each later phase [3, Sec.~3.5.1]. If
  a class changes, the plans are updated to match [1, Sec.~2.2.8, p.~19].
\end{itemize}

% ===========================================================================
\appendix
\renewcommand{\thesection}{\Alph{section}}

\section{Open Items}
The values below are not yet known.
\begin{description}
\tbditem{ref}{Project name and document reference.}
\tbditem{safety}{Safety-critical designation of each software item, from the
  hazard analysis.}
\end{description}

\clearpage
\section{Requirements Mapping Matrix}
Every requirement of NPR 7150.2D marked for an item's class is complied with
[1, App.~C, pp.~56--78], except those tailored out below
[1, Sec.~3.1.12, p.~22]. The tailoring was decided by the author on 10 October
2026 and is revisited at each later phase.

\begin{table}[H]
\centering
\small
\begin{tabular}{@{}lll@{}}
\toprule
SWE & Requirement & Reason \\
\midrule
141 & Software IV\&V for NASA project categories & NOT NASA \\
143 & Architecture review for NASA project categories & NOT NASA \\
174 & Planning data to the Center measurement repository & NOT NASA \\
148 & Reuse candidates to NASA's catalogue & NOT NASA \\
094 & Measurement data to the sponsoring organisation & NOT NASA \\
039 & Status reports from the software developer & NO SUPPLIER \\
040 & Products and records from the developer to NASA & NO SUPPLIER \\
042 & NASA access to the developer's source code & NO SUPPLIER \\
046 & Software schedule from the developer & NO SUPPLIER \\
206 & NASA access to the developer's models and simulations & NO SUPPLIER \\
045 & Joint NASA/developer audits & NO SUPPLIER \\
131 & IV\&V plan & NO IV\&V \\
178 & IV\&V access to development artifacts & NO IV\&V \\
179 & Responses to IV\&V issues and risks & NO IV\&V \\
032 & CMMI-DEV rated organisation & NO RATING \\
022 & Software assurance and safety plan per NASA-STD-8739.8 & NO ASSURANCE \\
205 & Safety-critical determination with the SMA organisation & NO ASSURANCE \\
015 & Two cost estimate models & NO BUDGET \\
151 & Conditions on the cost estimate & NO BUDGET \\
070 & Accredited models and simulations for flight qualification & BENCH ONLY \\
\bottomrule
\end{tabular}
\end{table}

\noindent\textbf{Legend: reason}
\begin{description}[leftmargin=3.9cm]
\item[\mbox{NOT NASA}] The System is not a NASA program or project; there is no NASA
  Center, catalogue or sponsoring organisation.
\item[\mbox{NO SUPPLIER}] The engineer develops all software; there is no separate
  developer or supplier.
\item[\mbox{NO IV\&V}] Independent verification and validation is not performed.
\item[\mbox{NO RATING}] There is no organisation to hold a CMMI rating.
\item[\mbox{NO ASSURANCE}] There is no independent software assurance or SMA
  organisation (ConOps constraint C-3); the author decides alone.
\item[\mbox{NO BUDGET}] The project has no development budget to estimate.
\item[\mbox{BENCH ONLY}] The System is demonstrated on a bench, not qualified for
  flight (ConOps constraint C-4).
\end{description}

\end{document}
